\documentclass[a4paper]{ltjsarticle}
\title{英単語テスト 概要}
\date{}

\begin{document}
\maketitle

\section{このコンテンツについて}
この単語テストは二種類の系統があり，
一つは第一回から第十七回までがある「システム英単語」準拠のもの（ファイル名がwordtest(N).pdf, 以後シス単テストと呼ぶ）と，
第一回から第三回以降まで鋭意製作中の通称「神単語 (ISBN: 978-4-7890-1855-5)」準拠のもの（ファイル名がkamiwordtest(N).pdf, 以降神単テストと呼ぶ）がある。
シス単テストは受験期に作っていたものであり，一応TOEICを想定したものではない。
あとシス単テストのほうが奇妙な例文が多い。
神単テストにはシス単テストの例文を一部流用しているがご容赦いただきたい。

クオートしているものを除いて，シス単テストは大抵の例文が，神単テストは注意書きのないもの以外のすべての例文の著作権は私にあるが，どのような範囲でも自由に利用して構わない。
間違い・誤植についてはいつでも受け付けている。

\section{解答方法}
紙でもらっているなら，そのまま書けば良い。
電子データでもらっているなら，PDFに直接書き込むか，適当な紙に答えだけ書いて写真などで提出すれば良い。
どちらでも採点の作業量は変わらないので，好きな方で。

\section{利用者}
神単テストの利用者情報 (2025/10/13)

\begin{tabular}{rl}
  Active: & 3名\\
  Inactive: & 1名\\
  No Data: & 1名
\end{tabular}



\end{document}